\chapter{Несоответствие: гидрид больше металла}\label{ch:misfit}

\section{Минимум теории упругости}

\paragraph{Деформация.} Если точка тела $\mathbf x$ сместилась на $\mathbf u(\mathbf x)$, то
малая деформация --- симметричная часть градиента смещения:
\begin{equation}
\eps_{ij}=\tfrac12\left(\partial_i u_j+\partial_j u_i\right).
\end{equation}
Диагональные компоненты --- удлинения вдоль осей, внедиагональные --- сдвиги.

\paragraph{Напряжение и закон Гука.} Напряжение $\sigma_{ij}$ --- сила на единицу площади
площадки с нормалью $i$ в направлении $j$. В изотропном упругом теле
\begin{equation}
\sigma=\lambda\,(\mathrm{tr}\,\eps)\,I+2\mu\,\eps\qquad\Longleftrightarrow\qquad\sigma=C:\eps,
\end{equation}
где $\lambda$, $\mu$ --- постоянные Ламе, $C$ --- тензор упругих модулей. В модели $E=90$~ГПа,
$\nu=0.34$ (цирконий при температуре выпадения; модули гидрида считаются такими же ---
по Puls 2012, разд.~2.2.3, у $\delta$-гидрида $E\approx80$--$85$~ГПа).

\paragraph{Равновесие.} Без внешних объёмных сил $\nabla\cdot\sigma=0$, то есть
$\partial_j\sigma_{ij}=0$ для каждого $i$.

\begin{exercise}
Найдите $\lambda=\frac{E\nu}{(1+\nu)(1-2\nu)}$ и $\mu=\frac{E}{2(1+\nu)}$ для $E=90$~ГПа, $\nu=0.34$.
Какая деформация отвечает напряжению 350~МПа (предел текучести)?
\end{exercise}

\section{Собственная деформация: мысленный опыт Эшелби}

Как описать включение, которое «хочет» быть больше, чем место, где оно сидит? Эшелби предложил
мысленный опыт в четыре шага (рис.~\ref{fig:F05_Eshelby}).

\fig{F05_Eshelby}{Мысленный опыт Эшелби. Кусок металла вырезают, он свободно превращается в
гидрид и меняет форму на $\eps^*$; его сжимают обратно до формы выреза, вставляют и отпускают.
Металл вокруг напрягается, включение остаётся сжатым.}

Превращение задаётся \emph{собственной деформацией} $\eps^*(\mathbf x)$ --- деформацией, которую
материал принял бы без соседей. Напряжение создаёт только \emph{разница} между настоящей
деформацией и собственной:
\begin{keyf}[Упругость с собственной деформацией]
\[
\sigma=C:\left(\eps-\eps^*\right),\qquad \nabla\cdot\sigma=0,\qquad \eps=\operatorname{sym}\nabla\mathbf u .
\]
\end{keyf}
У $\delta$-гидрида в осях пластинки (нормаль $n$, вдоль пластинки $t$, ось трубы $z$):
\begin{equation}
\eps^*_{nn}=\eps_n=0.072,\qquad \eps^*_{tt}=\eps^*_{zz}=\eps_t=0.0458
\end{equation}
(Carpenter 1973; ориентационное соотношение $\{111\}_\delta\parallel(0001)_\alpha$). Пластинка
расширяется во все стороны, но по нормали сильнее.

\begin{exercise}
Одномерная задача, которая объясняет многое. Периодический стержень длины $L$, на отрезке
длины $fL$ собственная деформация $\eps^*$, концы ячейки свободно смещаются (средняя сила
равна нулю). Из равновесия $d\sigma/dx=0$ найдите $\sigma(x)$. Почему в одномерии внутренних
напряжений нет вовсе, а в двух измерениях --- есть?
\end{exercise}

\section{Работа нагрузки и выгода $g$}

Пусть в металле уже есть напряжение $\sigma$ --- от нагрузки или от соседних пластинок. Когда в
нём появляется включение объёмом $V$ с несоответствием $\eps^*$, это напряжение совершает над
ним работу $V\,\sigma:\eps^*$ (здесь $\sigma:\eps^*=\sum_{ij}\sigma_{ij}\eps^*_{ij}$). Растяжение
помогает расширяющемуся включению, сжатие мешает. В модели эта работа на единицу объёма
нормирована на $\eps_n$, чтобы получить «эффективное напряжение»:

\begin{keyf}[Выгода зарождения от напряжений]
\[
g=\frac{\sigma:\eps^*_{\text{нов}}}{\eps_n}\ \ [\text{МПа}],\qquad\text{вклад в движущую силу: }\ \eps_n\,g .
\]
\end{keyf}

\fig{F06_Work}{Окружное напряжение $\sigma_\theta$ и две пластинки. У радиальной оно тянет вдоль
нормали, где несоответствие большое (7.2\%); у окружной --- вдоль пластинки, где оно меньше (4.6\%).
Отсюда упругий выбор ориентации нагрузкой.}

Для пластинки под углом $\psi$ к окружности при одном окружном напряжении:
\begin{equation}
g(\psi)=\sigma_\theta\,\frac{\eps_n\sin^2\psi+\eps_t\cos^2\psi}{\eps_n},\qquad
g_{\text{рад}}-g_{\text{окр}}=\frac{\eps_n-\eps_t}{\eps_n}\,\sigma_\theta=0.36\,\sigma_\theta .
\label{eq:gpsi}
\end{equation}

\begin{exercise}
Выведите формулу \eqref{eq:gpsi}: запишите $\eps^*$ пластинки в осях $(\theta,r)$ через поворот
на угол $\psi$. Посчитайте упругую разницу $g_{\text{рад}}-g_{\text{окр}}$ при 150~МПа и
переведите её в градусы ($\eps_n\Delta g$, затем 5.4~МПа/$^\circ$C). Сравните с измеренными
$0.08\,^\circ$C/МПа $\times150$~МПа. Во сколько раз упругость недооценивает эффект?
\end{exercise}

\begin{warnbox}
Упругая формула даёт в 15--20 раз меньше, чем измерено на синхротроне (Vizcaíno 2014). Часть
разницы объясняет пластичность (гл.~\ref{ch:plasticity}: втрое), остальное --- то, что
превращение идёт как сдвиг по плоскости габитуса, и расширение в плоскости пластинки снимается
течением. Поэтому в кинетической модели влияние нагрузки взято из опыта, а не из формулы.
\end{warnbox}

\section{Поле одной пластинки}

Решим задачу Эшелби для пластинки $5\times0.6$~мкм и посмотрим на $g$ для соседа той же
ориентации (рис.~\ref{fig:F07_Field}).

\fig{F07_Field}{Выгода $g$ зародыша той же ориентации от поля одной пластинки (упругость).
Перед кончиками --- красные лепестки: металл там растянут вдоль нормали, следующая пластинка
выгодна и продолжает цепочку. Над и под гранью --- синие области: там металл сжат.}

\begin{idea}
Пластинка «тянет» соседей к своим кончикам и отталкивает от граней. Поэтому гидриды не
рассыпаются, а выстраиваются в цепочки и стопки --- это \emph{автокаталитическое}
(самоускоренное) зарождение, о котором пишет Пульс. Направление цепочки задаёт ориентация
пластинок: окружные пластинки дают окружные цепочки.
\end{idea}

\section{Как поле считается: преобразование Фурье}

На сетке из сотен тысяч клеток с тысячами пластинок решать уравнения равновесия «в лоб» долго.
Помогает приём Хачатуряна (микроупругость, как в моделях фазового поля), рис.~\ref{fig:F08_Fourier}.

\fig{F08_Fourier}{Микроупругость через Фурье: собственная деформация на сетке $\to$ быстрое
преобразование Фурье $\to$ для каждой частоты $\mathbf k$ уравнение равновесия становится
системой $2\times2$ $\to$ обратное преобразование даёт напряжения.}

Подставим $\sigma=C:(\eps-\eps^*)$ в $\partial_j\sigma_{ij}=0$:
\[
C_{ijkl}\,\partial_j\partial_l u_k=C_{ijkl}\,\partial_j\eps^*_{kl}.
\]
В Фурье-пространстве производная --- умножение: $\partial_j\to ik_j$. Получаем для каждого
волнового вектора $\mathbf k$
\begin{equation}
K_{ik}(\mathbf k)\,\hat u_k=-i\,C_{ijkl}\,k_j\,\hat\eps^*_{kl},\qquad K_{ik}=C_{ijkl}k_jk_l .
\end{equation}
Матрица $K$ («акустический тензор») $2\times2$, и для изотропного тела её обратная известна
в явном виде:
\begin{equation}
K^{-1}=\frac{1}{\mu k^2}\left(I-\frac{\lambda+\mu}{\lambda+2\mu}\,\hat{\mathbf k}\hat{\mathbf k}^{\mathsf T}\right).
\end{equation}
Нашли $\hat{\mathbf u}$ --- нашли $\hat\eps$, обратным преобразованием --- $\eps(\mathbf x)$ и
$\sigma(\mathbf x)$. Стоимость --- два быстрых преобразования Фурье, $O(N\log N)$.

\paragraph{Цена: периодичность и среднее.} Фурье на сетке видит ячейку повторённой во все
стороны. Нулевая частота $\mathbf k=0$ из уравнения не определяется --- её задаём сами: среднюю
деформацию ячейки. Физически правильно: ячейка свободно расширяется, среднее напряжение от
гидридов равно нулю (поверхности трубы свободны).

\begin{warnbox}[Найдено в этой сессии]
В двумерном автомате долгое время стояла плоская деформация: $\eps_{zz}=0$. Тогда осевое
несоответствие гидридов ($\eps_t$ вдоль оси) некуда деть, и ячейка получает однородное сжатие
$-\lambda\langle\eps^*_{zz}\rangle$ в плоскости и $-(\lambda+2\mu)\langle\eps^*_{zz}\rangle$ по
оси. Сверка с МКЭ (CalculiX) показала: это сдвигает $g$ на $\approx-94$~МПа на каждый процент
объёма гидрида. Исправление --- обобщённая плоская деформация (средняя $\eps_{zz}$ равна средней
$\eps^*_{zz}$, опция \code{free\_z}). После него упругое поле автомата совпадает с МКЭ до 0.5~МПа.
\end{warnbox}

\begin{exercise}
Оцените сдвиг при плоской деформации для 178 ppm ($f\approx1.14\%$): посчитайте
$\langle\sigma_{xx}\rangle=\langle\sigma_{yy}\rangle=-\lambda f\eps_t$ и
$\langle\sigma_{zz}\rangle=-(\lambda+2\mu)f\eps_t$, затем
$g=(\eps_t\sigma_{xx}+\eps_n\sigma_{yy}+\eps_t\sigma_{zz})/\eps_n$. Сколько это в градусах?
\end{exercise}

\begin{codebox}
\textbf{Где в коде.} \code{ca\_hydride.py}: класс \code{Elastic} (обратная матрица $K^{-1}$,
напряжения по собственной деформации, опция \code{free\_z}), \code{plate\_eigen} (доля пластинки
в клетке и её $\eps^*$ в глобальных осях), \code{eigen\_components\_for}; проверка ---
\code{fe/stack\_auto.py}.
\end{codebox}
